%% kapa: exam=1 round=30 year=2019 period=1 form=A subject=economics subject_ko=경제학원론 seq=29 number=69
%% source: Q-Net 감정평가사 기출문제 (cst00309) artlSeq=5205205 file=extracted/1008/5205205/2019년 제30회 감정평가사 1차시험 문제지/2019년 제30회 감정평가사 1차시험 문제지/2019년도 제30회 감정평가사 1차 1교시 A형.hwp
%% conversion: kordoc 4.24.51 -> tools/gen1.py (auto); review=proofread-sonnet; flags=none
\begin{question}{69}
甲국과 乙국의 실질이자율과 인플레이션율은 다음 표와 같다. 명목이자소득에 대해 각각 25 \%의 세금이 부과될 경우, 甲국과 乙국의 세후 실질이자율은 각각 얼마인가? (단, 피셔효과가 성립한다.) \par
{\small\centering
\begin{tabular}{|c|c|c|}
\hline
 & \makecell{甲국} & \makecell{乙국} \\
\hline
\makecell{실질이자율} & \makecell{4 \%} & \makecell{4 \%} \\
\hline
\makecell{인플레이션율} & \makecell{0 \%} & \makecell{8 \%} \\
\hline
\end{tabular}\par}

\begin{choices}
  \item 3 \%, 1 \%
  \item 3 \%, 3 \%
  \item 3 \%, 9 \%
  \item 4 \%, 4 \%
  \item 4 \%, 12 \%
\end{choices}
\end{question}
\answer{1}
